

\subsection{Known Dichotomies and Hardness Results}\label{subsec: hardness results}

%This subsection introduces the notation and preliminary results used throughout the paper. Proofs of selected results are deferred to \autoref{app: disequality reduction}.

\paragraph{Dichotomy for Counting Eulerian Orientations}
$\newline$
  
For $\bowtie \,\in\{=,\ge,\le,>,<\}$, define
\(
\mathrm{HW}^{\bowtie}
=\left\{\sigma\in\{0,1\}^*\mid 
  \operatorname{wt}(\sigma)\mathrel{\bowtie}\frac{|\sigma|}{2}\right\}.
\)
With a slightly abuse of notation,
we also write $f\in\mathrm{HW}^{\bowtie}$ when
$\mathscr S(f)\subseteq\mathrm{HW}^{\bowtie}$.
An arity-$2d$ signature
$f$ is an \emph{$\eo$
signature} if $f\in\eoe$.
For a set $\mathcal F$ of $\eo$
signatures, we use
\[
\ceo(\mathcal F):=\holant{\neq_2}{\mathcal F}.
\]
For completeness, we now give the $\#\eo$ support notation used
in~\cite{GSS26}.  
For an $\eo$ signature $f$, and for all
$\beta,\gamma,\delta\in\mathscr S(f)$, define
\begin{align*}
f\in\POLMID
&\iff \beta\oplus\gamma\oplus\delta\in\mathscr S(f),\\
f\in\POLUP
&\iff \beta\oplus\gamma\oplus\delta
      \in\mathscr S(f)\cup\mathrm{HW}^{>},\\
f\in\POLDOWN
&\iff \beta\oplus\gamma\oplus\delta
      \in\mathscr S(f)\cup\mathrm{HW}^{<}.
\end{align*}
The latter two classes are the up- and down-quasi-polymorphism classes.
A \emph{perfect pairing} $M$ of the variables of an arity-$2d$
signature is a partition into $d$ unordered pairs.  It determines
\[
\{0,1\}^{M}
=\{x\in\{0,1\}^{2d}\mid x_p\ne x_q
  \text{ for every }\{p,q\}\in M\}.
\]
For any signature class $\mathscr C$, define
\[
\eo^{\mathscr C}
=\left\{f\mid f\text{ is an $\eo$ signature and }
  f|_{\{0,1\}^{M}}\in\mathscr C
  \text{ for every perfect pairing }M \text{ of } f\right\}.
\]





\begin{theorem}[Dichotomy for \#EO~\cite{MengWX25},\cite{GSS26}]
\label{thm: dichotomy for EO}
Let $\mathcal F$ be a finite set of complex-valued $\eo$ signatures.  If
\[
\bigl(\mathcal F\subseteq\POLUP
      \ \text{or}\ \mathcal F\subseteq\POLDOWN\bigr)
\quad\text{and}\quad
\bigl(\mathcal F\subseteq\eo^{\mathscr A}
      \ \text{or}\ \mathcal F\subseteq\eo^{\mathscr P}\bigr),
\]
then $\ceo(\mathcal F)$ is computable in polynomial time.  Otherwise,
$\ceo(\mathcal F)$ is \#P-hard.
\end{theorem}
The two alternatives are uniform over $\mathcal F$: signatures from the
up and down classes cannot be mixed, nor can
$\eo^{\mathscr A}$ and $\eo^{\mathscr P}$ be mixed.


\begin{lemma}\label{lm: conjugation_polymorphism_affine support}
    Suppose $\mathcal{F}$ is a conjugate closed signature set.
    If $\widehat{\mathcal{F}}\subseteq \eog$ or $\widehat{\mathcal{F}}\subseteq \eol$, then every $\widehat{f}\in \widehat{\mathcal{F}}$ is $\eo$.
    Moreover, if $\widehat{\mathcal{F}}\subseteq \POLUP$ or $\widehat{\mathcal{F}}\subseteq \POLDOWN$, then every $\widehat{f}\in \widehat{\mathcal{F}}$ has affine support.
\end{lemma}

\begin{proof}
    We assume $\widehat{\mathcal{F}}\subseteq \eog$, the other case can be handled similarly.
    By \cref{lm: conjugate closed and ars closed}, $\widehat{\mathcal{F}}$ is arrow reversal closed.
    Hence $\widehat{\mathcal{F}}\subseteq \eog\cap \eol$, i.e., $\widehat{\mathcal{F}}$ is a set of $\eo$ signatures.
    We assume $\widehat{\mathcal{F}}\subseteq \POLUP$ without loss of generality.
    Pick arbitrary $\widehat{f}\in \widehat{\mathcal{F}}$ and $\alpha,\beta,\gamma\in \mathscr{S}(\widehat{f})$, we have $\alpha\oplus \beta\oplus \gamma\in \mathscr{S}(\widehat{f})\cup \eosg$.
    By \cref{lm: conjugate closed and ars closed}, $\widehat{f}^{\textsc ar}\in \widehat{\mathcal{F}}$.
    Notice that $\overline{\alpha},\overline{\beta},\overline{\gamma}\in \mathscr{S}(\widehat{f}^{\textsc ar})$.
    Then $\overline{\alpha\oplus\beta\oplus\gamma}\in \mathscr{S}(\widehat{f}^{\textsc ar})\cup \eosg$.
    It follows that $\alpha\oplus\beta\oplus\gamma\in \mathscr{S}(\widehat{f})\cup \eosl$.
    Thus, $\alpha\oplus\beta\oplus\gamma\in (\mathscr{S}(\widehat{f})\cup \eosl)\cap (\mathscr{S}(\widehat{f})\cup \eosg)=\mathscr{S}(\widehat{f})$.
    Therefore, $\widehat{f}$ has affine support.
\end{proof}


\paragraph{Dichotomy for Holant with an Odd Arity Signature}
$\newline$
A signature of arity $k$ is \emph{single-weighted} if its support is
contained in the strings of one Hamming weight $d$.  In that case define
the associated $\eo$ signature
\[
f_{\to\eo}=
\begin{cases}
f\otimes\Delta_0^{\otimes(2d-k)},&2d\ge k,\\
f\otimes\Delta_1^{\otimes(k-2d)},&2d<k.
\end{cases}
\]
Also set $f|_{\eo}:=f|_{\mathrm{HW}^{=}}$, and extend both operations
elementwise to signature sets:
\[
\mathcal F|_{\eo}=\{f|_{\eo}:f\in\mathcal F\},\qquad
\mathcal F_{\to\eo}=\{f_{\to\eo}:f\in\mathcal F\}.
\]
Finally,
\[
\mathcal M
=\{f:f(\sigma)=0\text{ whenever }\operatorname{wt}(\sigma)>1\}
\]
is the class of weighted-matching signatures.
Notice that 
\[
X\mathcal M
=\{f:f(\sigma)=0\text{ whenever }\operatorname{wt}(\sigma)<n-1\}
\]


\begin{theorem}[Dichotomy for $\holodd$ \cite{Holant_Odd}]\label{thm: holant odd}
    Let $\mathcal{F}$ be a set of signatures that contains a non-zero signature of odd arity, then $\Holant(\mathcal{F})$ is \#P-hard unless:
    \begin{enumerate}
        \item $\widehat{\mathcal{F}}$ only contains $\eog$ signatures (or $\eol$ signatures respectively), and $\widehat{\mathcal{F}}|_{\eo}\subseteq \POLUP$ or $\widehat{\mathcal{F}}|_{\eo}\subseteq \POLDOWN$, and $\widehat{\mathcal{F}}|_{\eo}\subseteq \eo^\mathscr{A}$ or $\widehat{\mathcal{F}}|_{\eo}\subseteq \eo^\mathscr{P}$;
        
        \item 
        All signatures in $\widehat{\mathcal{F}}$ are single-weighted. 
        There exists a signature in $\eosg$ and another in $\eosl$ belonging to $\widehat{\mathcal{F}}$.  
        In addition, $\widehat{\mathcal{F}}_{\to\eo}\subseteq \POLUP$ or $\widehat{\mathcal{F}}_{\to\eo}\subseteq \POLDOWN$, and $\widehat{\mathcal{F}}_{\to\eo}\subseteq \eo^\mathscr{A}$ or $\widehat{\mathcal{F}}_{\to\eo}\subseteq \eo^\mathscr{P}$;

        

        \item $\widehat{\mathcal{F}}\subseteq\langle \mathcal{M}\rangle$ or $\widehat{\mathcal{F}}\subseteq\langle X\mathcal{M}\rangle$;

        \item $\mathcal{F}\subseteq\mathscr{T}$;
        
        \item $\mathcal{F}$ is $\mathscr{A}$-transformable;

        \item $\mathcal{F}$ is $\mathscr{P}$-transformable;

        \item $\mathcal{F}$ is $\mathscr{L}$-transformable;
    \end{enumerate}
    We denote Case 1-7 as condition $(\mathcal{PC})$.
\end{theorem}

\begin{theorem}[\cite{Holant_Odd}]\label{thm: decomposition lemma}
    Let $\mathcal{F}$ be a set of signatures and $f$, $g$ be two signatures. Then

$$\Holant(f,g,\mathcal{F})\equiv_T\Holant(f\otimes g,\mathcal{F})$$ 
holds unless $\widehat{\mathcal{F}}'=\widehat{\mathcal{F}}\cup\{\widehat{f}\otimes\widehat{g}\}$ only contains $\eog$ signatures (or $\eol$ signatures respectively).
\end{theorem}




\begin{lemma}[Decomposition lemma]\label{lm: decomposition}
    Let $\mathcal{F}$ be a conjugate closed set of signatures.
    If $f\in \mathcal{F}$ has a factorization $g\otimes h$, then $\Holant(g, h, \mathcal{F})\equiv_T\Holant( \mathcal{F})$.
    %In the setting of $\holant{\neq_2}{\hF}$, $\holant{\neq_2}{\widehat{g}, \widehat{h}, \hF}\equiv_T\holant{\neq_2}{\hF}$.
\end{lemma}

\begin{proof}
    By \cref{thm: decomposition lemma}, we only need to consider the case that $\widehat{\mathcal{F}}$ only contains $\eog$ signatures (or $\eol$ signatures respectively).
    Since $\mathcal{F}$ is conjugate closed, $\widehat{\mathcal{F}}$ is $\eo$ by \cref{lm: conjugation_polymorphism_affine support}.
    By \cref{thm: dichotomy for EO}, either $\Holant(\mathcal{F})\equiv_T\holant{\neq_2}{\widehat{\mathcal{F}}}$ is \#P-hard, then trivially $\Holant(g,h,\mathcal{F})\equiv_T\Holant(\mathcal{F})$, or $\holant{\neq_2}{\hF}$ is tractable.
    In the tractable case, $\widehat{\mathcal{F}}\subseteq \POLUP$ or $\widehat{\mathcal{F}}\subseteq \POLDOWN$ and $\widehat{\mathcal{F}}\subseteq \eo^{\mathscr{A}}$ or $\widehat{\mathcal{F}}\subseteq \eo^{\mathscr{P}}$.
    By \cref{lm: conjugation_polymorphism_affine support}, every $\widehat{f}\in \mathcal{F}$ has affine support.
    Combined with $\widehat{\mathcal{F}}\subseteq \eo^{\mathscr{A}}$ or $\widehat{\mathcal{F}}\subseteq \eo^{\mathscr{P}}$, we have $\widehat{\mathcal{F}}\subseteq \mathscr{A}$ or $\widehat{\mathcal{F}}\subseteq\mathscr{P}$.
    Since $\widehat{f}=\widehat{g}\otimes\widehat{h}$, $\widehat{g}$ and $\widehat{h}$ are both affine or both product.
    Then trivially $\holant{\neq_2}{\widehat{g},\widehat{h},\widehat{\mathcal{F}}}\equiv_T \holant{\neq_2}{\widehat{\mathcal{F}}}$, which is equivalent to $\Holant(g, h, \mathcal{F})\equiv_T\Holant( \mathcal{F}).$ 
\end{proof}



\begin{lemma}\label{lm: odd holant with conjugation}
    Let $\mathcal{F}$ be a conjugate closed set of signatures containing a nonzero signature of odd arity.
    If $\mathcal{F}$ does not satisfy condition \rm{(\ref{cond: tractable classes})},
    then $\Holant(\mathcal{F})$ is \#P-hard.
\end{lemma}


\begin{proof}
    Notice that Case 4-7 in $(\mathcal{PC})$ is {\rm(\ref{cond: tractable classes})}.
    By \cref{thm: holant odd}, we only need to show that Case 1-3 in $(\mathcal{PC})$ cannot happen.
    \begin{enumerate}
        \item 
        Case 1 in $(\mathcal{PC})$ happens.
        By \cref{lm: conjugation_polymorphism_affine support}, $\hF$ is $\eo$ and has affine support.
        Combined with $\widehat{\mathcal{F}}|_{\eo}\subseteq \eo^\mathscr{A}$ or $\widehat{\mathcal{F}}|_{\eo}\subseteq \eo^\mathscr{P}$, we have $\hF\subseteq \mathscr{A}$ or $\hF\subseteq \mathscr{P}$.
        So $\F$ is $\mathscr{A}$-transformable or $\mathscr{P}$-transformable, contradicting that $\F$ does not satisfy \rm{(\ref{cond: tractable classes})}.

        \item 
        Case 2 in $(\mathcal{PC})$ happens.
        By \cref{lm: conjugation_polymorphism_affine support}, $\widehat{\mathcal{F}}_{\to\eo}$ has affine support.
        Combined with $\widehat{\mathcal{F}}_{\to\eo}\subseteq \eo^\mathscr{A}$ or $\widehat{\mathcal{F}}_{\to\eo}\subseteq \eo^\mathscr{P}$, we have
        $\widehat{\mathcal{F}}_{\to\eo}\subseteq \mathscr{A}$ or  $\widehat{\mathcal{F}}_{\to\eo}\subseteq \mathscr{P}$.
        So $\hF\subseteq \mathscr{A}$ or $\hF\subseteq \mathscr{P}$,
        contradicting \rm{(\ref{cond: tractable classes})}.

        \item 
        Case 3 in $(\mathcal{PC})$ happens.
        Pick an arbitrary $\widehat{f}\in \hF$.
        Write its unique prime factorization $\widehat{f}=\widehat{q_1}\otimes \widehat{q_2}\otimes\cdots\otimes \widehat{q_s}.$
        We first assume $\hF\in \braket{\widehat{\mathcal{M}}}.$
        Then every $\widehat{q_i}$ is supported on Hamming weight at most 1.
        By \cref{lm: conjugate closed and ars closed}, $\hF$ is arrow reversal closed.
        Taking {\sc ars} doesn't change factorization and irreducibility.
        So $\widehat{f}^{\textsc ar}=\widehat{q_1}^{\textsc ar}\otimes \widehat{q_2}^{\textsc ar}\otimes\cdots\otimes \widehat{q_3}^{\textsc ar}\in \hF$, where every $\widehat{q_i}^{\textsc ar}$ is also supported on Hamming weight at most 1.
        This is impossible if some $\widehat{q_i}$ has arity greater than 2.
        Thus, every $q_i$ has arity at most 2.
        So $\widehat{f}\in {\mathscr{T}}$, and $\hF\subseteq \mathscr{T}$.
        After a $K$-transformation,  $\F\subseteq \mathscr{T}$.
        Contradiction.
    \end{enumerate}
\end{proof}


\paragraph{Dichotomy for Quaternary Signatures}
$\newline$
%Let $f$ be an arity 4 signature supported on even Hamming weights.
%The \emph{eight-vertex model} is the problem $\Holant(\neq_2\mid f).$
%Moreover, if $f(0000)=f(1111)=0$, then $\Holant(\neq_2\mid f)$ is a \emph{six-vertex model}.

\begin{theorem}[Dichotomy for six-vertex models]\label{thm: six vertex model}
Let $f$ be a 4-ary signature with the signature matrix
$M_{x_1x_2, x_4x_3}=\begin{bmatrix}
0 & 0 & 0 & a\\
0 & b & c & 0\\
0 & z & y & 0\\
x & 0 & 0 & 0\\
\end{bmatrix}$, then
$\Holant(\neq_2\mid f)$ is \#P-hard except for the following cases:
\begin{itemize}
\item $f\in\mathscr{P}$;
\item $f\in\mathscr{A}$;
\item there is a zero in each pair $(a,x), (b,y), (c,z)$;
\end{itemize}
in which cases Holant$(\neq_2\mid f)$ is computable in polynomial time.
For the last tractable case, $f\in \braket{\mathcal{M}}$ or $f\in \braket{X\mathcal{M}}$. 
\end{theorem}

For an arity 4 signature $h$, we say that $h$ has \emph{$H$-form}
if $h\in\mathrm{HW}^{\ge}$ or $h\in\mathrm{HW}^{\le}$, and define its
reduced form by $r(h):=h|_{\eo}$.  Equivalently, $M(h)$ has one of the
following two forms:
\[
\begin{bmatrix}
0&0&0&b\\
0&c&d&\ast\\
0&w&z&\ast\\
y&\ast&\ast&\ast
\end{bmatrix}
\quad\text{or}\quad
\begin{bmatrix}
\ast&\ast&\ast&b\\
\ast&c&d&0\\
\ast&w&z&0\\
y&0&0&0
\end{bmatrix},
\qquad
M(r(h))=
\begin{bmatrix}
0&0&0&b\\
0&c&d&0\\
0&w&z&0\\
y&0&0&0
\end{bmatrix}.
\]
If every right-hand vertex of an instance of
$\holant{\neq_2}{h}$ is assigned $h$, then the total Hamming weight on
the right is exactly half the number of incident edges.  Hence, when
$h$ has $H$-form, every local assignment occurring in a nonzero global
assignment has weight two.  Therefore $\holant{\neq_2}{h}\equiv_T\holant{\neq_2}{r(h)}.$

Thus the reduced problem is a six-vertex model, equivalently a quaternary
$\#\mathrm{EO}$ problem.  We place it together with the following
eight-vertex result since both serve as quaternary vertex-model base
cases used later.

\iffalse
\begin{theorem}[\cite{eight-vertex_model}]\label{thm: eight-vertex model}
    Let $\widehat{f}$ be a non-zero signature with $M(\widehat{f})=\left[\begin{smallmatrix}
    a & 0 & 0 & b\\
    0 & c & d & 0\\
    0 & d & c & 0\\
    b' & 0 & 0 & \overline{a}\\
    \end{smallmatrix}\right]$ and $\hF$ be a set of signatures.
    Then, $\holant{\neq_2}{\widehat f,\hF}\equiv_T \holant{\neq_2}{\widehat g, \hF}$ where $M(\widehat{g})=\left[\begin{smallmatrix}
    |a| & 0 & 0 & b\\
    0 & c & d & 0\\
    0 & d & c & 0\\
    b' & 0 & 0 & |a|\\
    \end{smallmatrix}\right]$, and $\holant{\neq_2}{\widehat g}$ is \#P-hard unless $\widehat{g}$ has support size 2, 4 or 8.
    Moreover, if $\widehat{g}$ has support size 8, then one of the following holds:
    \begin{enumerate}
        \item $\holant{\neq_2}{\widehat{g}}$ is \#P-hard.

        \item There exists $T\in\mathrm{GL}_2(\mathbb{C})$ such that $\#\mathrm{CSP}_2(T\hF)\le_T\holant{\neq_2}{\widehat{g},\hF}.$

        \item $b=b'$ and $|a|=|b|=|c|=|d|$.
    \end{enumerate}
\end{theorem}
\fi

Following~\cite{HuangFu4Regular}, let $\mathscr H$ denote the set of
arity-four signatures $f$ for which $\widehat f$ has $H$-form and
$\holant{\neq_2}{r(\widehat f)}$ is tractable.

\begin{theorem}[Dichotomy for a single quaternary signature
\cite{HuangFu4Regular}]\label{thm: single quaternary dichotomy}
    Let $f$ be an arity-four signature.  Then $\holant{=_2}{f}$ is
    \#P-hard unless one of the following holds:
    \begin{enumerate}
        \item $f\in\langle K\mathcal M\rangle$ or
        $f\in\langle KX\mathcal M\rangle$;
        \item $f$ is $\mathscr A$-transformable;
        \item $f$ is $\mathscr L$-transformable;
        \item $f$ is $\mathscr P$-transformable;
        \item $f\in\mathscr H$;
        \item $f\in\langle\mathscr T\rangle$.
    \end{enumerate}
    In these cases, the problem is computable in polynomial time.
    %Also, cases 2, 3, 4 are tractable cases for the eight vertex model $\Holant(\neq_2\mid \widehat{f})$.
\end{theorem}


\paragraph{Other Hardness Results}
The following is a standard lemma by polynomial interpolation. See for example, \cite{jcbook}.
\begin{restatable}{lemma}{InterpolationBinary}\label{lm: 2 by 2 interpolation}
    Let $f$ and $g$ be nonzero binary signatures with $M(f)=P^{-1}\left[\begin{smallmatrix}
    \lambda_1 & 0\\
    0 & \lambda_2\\
    \end{smallmatrix}\right]P$ and
    $M(g)=P^{-1}\left[\begin{smallmatrix}
    1 & 0\\
    0 & 0\\
    \end{smallmatrix}\right]P$ for some invertible matrix $P$.
    If $\lambda_1\neq 0$ and $\frac{\lambda_2}{\lambda_1}$ is not a root of unity,
    then $$\Holant(g, \mathcal{F})\leqslant_T\Holant(f, \mathcal{F})$$ for any signature set $\mathcal{F}$.
\end{restatable}

\begin{theorem}[\cite{real_holantc}]\label{thm: CSP2}
    Let $\mathcal{F}$ be any set of signatures. 
    Then $\#\operatorname{CSP}_2(\mathcal{F})$ is $\#\operatorname{P}$-hard unless $\mathcal{F}\subseteq\mathscr{A}$ or $\mathcal{F}\subseteq\alpha\mathscr{A}$ or $\mathcal{F}\subseteq\mathscr{P}$ or $\mathcal{F}\subseteq\mathscr{L}$, in which cases the problem is computable in polynomial time.
\end{theorem}

\begin{lemma}[\cite{jcbook}]\label{lm: CSP2 to holant =4}
    $\#\mathrm{CSP}_2(\mathcal{F})\leqslant_T\Holant(=_4, \mathcal F)$.
\end{lemma}

\begin{restatable}{lemma}{CSPwithEQ}\label{lm: CSP with =2}
    Suppose $\mathcal{F}$ doesn't satisfy \eqref{cond: tractable classes}.
    Then for every $T\in \mathrm{GL}_2(\C)$, $\#\mathrm{CSP}_2(T(\mathcal{F}\cup \{=_2\}))$ is \#P-hard.
\end{restatable}

\begin{restatable}{lemma}{DisFourgivesHard}\label{lm: disequality 4 gives dichotomy}
    Let $\F$ be a set of signatures. Then $\holant{\neq_2}{\neq_4,\hF}$ is either polynomial-time solvable or $\#\operatorname{P}$-hard.
\end{restatable}


\begin{restatable}{lemma}{ConjuDisFourHard}\label{lem: In conjugate closed setting disequality 4 gives hardness}
    Let $\F$ be a conjugate closed signature set. If $\F$ does not satisfy the condition \rm{(\ref{cond: tractable classes})}, then $\holant{\neq_2}{\neq_4,\hF}$ is $\#\operatorname{P}$-hard.
\end{restatable}

The proofs of~\cref{lm: CSP with =2}, ~\cref{lm: disequality 4 gives dichotomy} and~\cref{lem: In conjugate closed setting disequality 4 gives hardness} are given in \autoref{app: disequality reduction}.
